% Part: normal-modal-logic
% Chapter: axioms-systems
% Section: introduction

\documentclass[../../../include/open-logic-section]{subfiles}

\begin{document}

\olfileid{nml}{prf}{int}\sourcecorrection{OLMD-022}{મૂળમાં આ પરિચયનું પ્રકરણ-ઓળખચિહ્ન axs છે, જ્યારે તેને આયાત કરતું પ્રકરણ prf ઓળખચિહ્ન વાપરે છે; અહીં બંને એકસરખાં કર્યાં છે.}

\olsection{પરિચય}

આપણી પાસે મૂળભૂત મોડલ ભાષા માટે મોડલ નિદર્શનો પર આધારિત અર્થવિજ્ઞાન છે.
!!a{formula}ની પ્રામાણ્યતા એટલે બધા નિદર્શનોનાં બધાં જગતોમાં તેની સત્યતા;
નિદર્શનોના કોઈ વર્ગ કે ફ્રેમને અનુલક્ષીને પ્રામાણ્યતા એટલે તે વર્ગનાં બધાં
નિદર્શનોનાં, અથવા તે ફ્રેમ પર આધારિત બધાં નિદર્શનોનાં, બધાં જગતોમાં તેની
સત્યતા. તર્કશાસ્ત્ર સામાન્ય રીતે પ્રામાણ્યતાનાં આવાં અર્થાનુસારી વર્ણનોને
!!{derivability}ના પ્રમાણસિદ્ધાંતલક્ષી ખ્યાલ સાથે જોડે છે. આપણો હેતુ
કોઈ તંત્રમાં !!{derivability} એવી રીતે વ્યાખ્યાયિત કરવાનો છે કે !!a{formula}
પ્રામાણ્ય હોય તો અને માત્ર તો જ તે !!{derivable} હોય.

સૌથી સરળ અને ઐતિહાસિક રીતે સૌથી જૂનાં !!{derivation} તંત્રો હિલ્બર્ટ-પ્રકારનાં,
અથવા સ્વયંસિદ્ધિમૂલક, !!{derivation} તંત્રો છે. અનેક મોડલ તર્કો માટે આવાં
!!{derivation} તંત્રો રચવાં પ્રમાણમાં સહેલાં છે; અધિસૈદ્ધાંતિક અભ્યાસમાં પણ તે સરળ છે
(દાખલા તરીકે, યથાર્થતા અને પૂર્ણતા સાબિત કરવા). જોકે, કુદરતી નિષ્કર્ષણનાં
તંત્રોની સરખામણીમાં તેમાં !!{formula}s સાબિત કરવાં ઘણાં અઘરાં છે.

હિલ્બર્ટ-પ્રકારનાં !!{derivation} તંત્રોમાં !!a{formula}ની !!a{derivation}
એ ચોક્કસ સ્વયંસિદ્ધિઓથી શરૂ થઈને થોડા અનુમાન-નિયમો વડે સંબંધિત !!{formula}
સુધી પહોંચતી !!{formula}sની શ્રેણી છે. !!{derivation} તંત્ર અર્થવિજ્ઞાન સાથે
મેળ ખાય તે માટે આપણે ખાતરી કરવી પડે કે !!{derivable} સૂત્રો બધાં
નિદર્શનોમાં સત્ય છે (અથવા બધાં સ્વયંસિદ્ધિઓ સત્ય હોય એવાં બધાં
નિદર્શનોમાં સત્ય છે). પહેલાં આપણે આ માટે જરૂરી મોડલ તર્કોના કેટલાક
ગુણધર્મો અલગ તારવીશું: આ ``સામાન્ય'' મોડલ તર્કો છે. સામાન્ય મોડલ તર્કો
માટે અનુમાનના માત્ર બે નિયમો ધારવા પડે છે: મોડસ પોનેન્સ અને આવશ્યકીકરણ.
સ્વયંસિદ્ધિઓ તરીકે આપણે પુનરુક્તિઓના બધા પ્રતિસ્થાપન દૃષ્ટાંતો અને,
જે મોડલ તર્કની ચર્ચા કરીએ તેના આધારે, કેટલીક મોડલ સ્વયંસિદ્ધિઓ લઈએ
છીએ. જો આપણને બધાં નિદર્શનોના વર્ગમાં જ રસ હોય, તોપણ~$\Ax{K}$\iftag{notprvDiamond}{}{ અને~$\Ax{Dual}$}ના
બધા પ્રતિસ્થાપન દૃષ્ટાંતોને સ્વયંસિદ્ધિઓ ગણવાં પડે છે. માત્ર આટલું
ન્યૂનતમ સામાન્ય મોડલ તર્ક~$\Log K$ પેદા કરે છે.

\begin{defn}
\emph{મોડસ પોનેન્સ}નો નિયમ આ અનુમાન-ઢાંચો છે:
\begin{prooftree}
\AxiomC{$!A$}
\AxiomC{$!A \lif !B$}
\RightLabel{\MP}
\BinaryInfC{$!B$}
\end{prooftree}
જો $!C \ident !A \lif !B$ હોય, તો આપણે કહીએ છીએ કે !!{formula}s $!A$, $!C$
માંથી મોડસ પોનેન્સ વડે !!a{formula}~$!B$ \emph{ફલિત થાય છે}.
\end{defn}

\begin{defn}
\emph{આવશ્યકીકરણ}નો નિયમ આ અનુમાન-ઢાંચો છે:
\begin{prooftree}
\AxiomC{$!A$}
\RightLabel{\Nec}
\UnaryInfC{$\Box !A$}
\end{prooftree}
જો $!B \ident \Box !A$ હોય, તો આપણે કહીએ છીએ કે !!{formula}~$!A$માંથી
આવશ્યકીકરણ વડે !!{formula}~$!B$ ફલિત થાય છે.
\end{defn}

\begin{defn}
સ્વયંસિદ્ધિઓના ગણ~$\Sigma$માંથી \emph{!!{derivation}} એ !!{formula}s
$!B_1$, $!B_2$, \dots, $!B_n$ની એવી શ્રેણી છે કે દરેક $!B_i$ આમાંથી
એક હોય:
\begin{enumerate}
\item પુનરુક્તિનો પ્રતિસ્થાપન દૃષ્ટાંત, અથવા
\item $\Sigma$માંના !!a{formula}નો પ્રતિસ્થાપન દૃષ્ટાંત, અથવા
\item $j$, $k < i$ હોય એવાં બે !!{formula}s $!B_j$, $!B_k$માંથી
  મોડસ પોનેન્સ વડે ફલિત થતું સૂત્ર, અથવા
\item $j < i$ હોય એવા !!a{formula}~$!B_j$માંથી આવશ્યકીકરણ વડે
  ફલિત થતું સૂત્ર.
\end{enumerate}
જો $!B_n \ident !A$ હોય એવી !!a{derivation} હોય, તો આપણે કહીએ છીએ કે
$!A$ એ~$\Sigma$માંથી \emph{!!{derivable}} છે; સંકેતમાં $\Sigma \Proves !A$.
\end{defn}

આ વ્યાખ્યા અનુસાર !!{derivable} સૂત્રોનો ગણ સામાન્ય મોડલ તર્ક બને છે,
અને દરેક !!{derivable} !!{formula} એવા દરેક નિદર્શનમાં સત્ય છે જેમાં
દરેક સ્વયંસિદ્ધિ સત્ય છે. !!{derivation}sના આ ગુણધર્મને \emph{યથાર્થતા}
કહે છે. તેના વિપરીત ગુણધર્મ, \emph{પૂર્ણતા}, સાબિત કરવો વધુ અઘરો છે.

\end{document}
